\begin{lemma}\label{lem:decic-318-cubic-fibre}
Put $F=\mathbb Q(\zeta_3)$ and $X=X_1(3,18)_F$.
An effective noncuspidal $F$-rational divisor of degree $5$ on $X$
cannot reduce at $7$ or $13$ to one of the cubic places listed
below together with a cuspidal divisor of degree $2$.
\end{lemma}
\begin{proof}
At $7$, the cubic indices are $4,5,7,8$. At $13$, they are
$15,19,28,29,31,32,35,36,37,38$. The indices refer to the
complete $27$-automorphism enumerations. Let $p$ denote the chosen prime.
Each $N_i$ is fixed by a nonidentity automorphism $h_i$ of order $3$.
These level changes preserve the nonzero point $9Q_{18}$ of order $2$.
They act freely on geometric noncuspidal points: an elliptic
curve automorphism of order $3$ modulo sign has $\deg(u-1)=3$
and cannot fix nonzero $2$-torsion.

Write $\pi_i:X\to Y_i=X/\langle h_i\rangle$.
The quotient is smooth at the chosen prime, since the action
has order prime to its residue characteristic, and $\pi_i$ is \'etale at $N_i$.
The closed cubic is therefore the full fibre of a point
$q_i\in Y_i(\mathbb F_p)$.
For every effective rational two-cusp divisor $C_2$,
restriction of
\[
 (h_i^*-1)H^0(X_{\mathbb F_p},\Omega^1)
                 \longrightarrow H^0(N_i+C_2,\Omega^1)
\]
has rank $4$. The source space has dimension $6$ or $8$.
The computation tests $684$ cases at $7$ and $1710$ at $13$,
including double cusps, using the full Jacobian differential space.
A fresh native replay checks every stabilizer and every differential
kernel independently.

Suppose that $\overline D=N_i+C_2$.
Lift $C_2$ cuspidally over $R=\mathbb Z_p$, and lift $q_i$ to
$Y_i(R)$ using smoothness. Put $\alpha=h_{i*}-1$.
The class
\[
 [h_iD-D-h_iC_2+C_2]
\]
is torsion, by rank zero over $F$ and Manin--Drinfeld, and has
zero reduction. The odd unramified formal group has no torsion,
so this class is zero.
Thus $D$ lies in the same $\alpha$-Abel fibre as
$C_2+\pi_i^*q_i$.

The family $C_2+\pi_i^*Y_i$ is a smooth closed formal family of
dimension $1$ in $X^{(5)}$ at this residue divisor.
Its moving support is disjoint from $C_2$, and the three moving
branches are \'etale over $Y_i$, so their coefficients recover
the parameter of $Y_i$. It lies in the $\alpha$-Abel fibre.
The rank-$4$ restriction gives four independent equations for
that fibre in the smooth five-dimensional symmetric power.
Their zero locus is formally smooth of dimension $1$.
The tangent map of the known family is an isomorphism onto this
zero locus, so the formal implicit function theorem identifies
them. Hence $D$ contains $C_2$, contradicting its noncuspidal support.
\end{proof}

\begin{lemma}\label{lem:decic-318-quartic-fibre}
Let $F=\mathbb Q(\zeta_3)$ and $X=X_1(3,18)_F$.
An effective noncuspidal $F$-rational divisor of degree $5$ on $X$
cannot reduce at $7$ to a noncuspidal quartic in either of the
orbits numbered $12$ and $17$ in the complete $27$-automorphism
enumeration, together with a rational cusp.
\end{lemma}
\begin{proof}
For a level structure $(E,P_3,Q_{18})$, let $T=9Q_{18}$ and
let $\phi:E\to E'=E/\langle T\rangle$.
Write $T'$ for the nonzero point in the kernel of the dual
$2$-isogeny. The formula
\[
 \tau(E,P_3,Q_{18})=
      (E',\phi P_3,2\phi Q_{18}+T')
\]
defines an $F$-rational involution of $X$.
Indeed, the new point of order $18$ is the sum of a point of
order $9$ and $T'$, and the Weil pairing is raised to the power
$2\cdot2\equiv1\pmod3$. Applying the formula twice gives
$(-P_3,-Q_{18})$, which is the original point of the modular curve.

The quartic representatives $N_{12}$ and $N_{17}$ are fixed
as closed places by $h_i=g_i^{-1}\tau g_i$, where
\[
 g_{12}(P,Q)=(P+12Q,Q),\qquad
 g_{17}(P,Q)=(P+12Q,P+Q).
\]
In both cases $h_i$ acts on the quartic residue field as
Frobenius squared. Thus the quotient $\pi_i:X\to Y_i=X/\langle h_i\rangle$
is \'etale at $N_i$, and
$N_i=\pi_i^*q_i$ for a quadratic place $q_i$ on the quotient.
The anti-invariant differential space has dimension $6$.
For every one of the $18$ rational cusps $c$, restriction of
\[
 (h_i^*-1)H^0(X_{\mathbb F_7},\Omega^1)
       \longrightarrow H^0(N_i+c,\Omega^1)
\]
has rank $3$. All $36$ restrictions, including their exact
Riemann--Roch kernels, are retained in the certificate.

Suppose that $\overline D=N_i+c$.
Let $R=\mathbb Z_7$, lift $c$ cuspidally, and lift $q_i$ to a
quadratic trace divisor $E_2$ on $Y_i/R$.
Put $\alpha=h_{i*}-1$ and $D_0=c+\pi_i^*E_2$.
Rank zero of $J_X(F)$, Manin--Drinfeld, and
$\alpha\pi_i^*=0$ show that $\alpha[D-D_0]$ is torsion.
It has zero reduction, so the odd unramified formal group
implies $\alpha[D-D_0]=0$.

The formal family $c+\pi_i^*Y_i^{(2)}$ through $D_0$ is smooth
of dimension $2$ and is a closed formal immersion into $X^{(5)}$.
The moving support is disjoint from $c$, and $\pi_i$ is \'etale
at the moving points, so the two moving coefficients can be
recovered from their pullbacks.
The rank-$3$ restriction gives three independent cotangent
equations for the $\alpha$-Abel fibre in the smooth
five-dimensional symmetric power. Their zero locus is smooth
of dimension $2$, and contains the preceding family with the
same tangent space. The formal implicit function theorem
identifies them. Consequently $c\leq D$, a contradiction.
\end{proof}

\begin{lemma}\label{lem:decic318-last-cuspidal-classes}
Let $F=\mathbb Q(\zeta_3)$ and $X=X_1(3,18)$ over $F$.
After the preceding finite and relative tests, every noncuspidal
$F$-rational effective divisor of degree $5$ on $X$, after applying
an automorphism in the order-$27$ marking group $G$, is linearly
equivalent to one of the following two cuspidal divisors:
\[
\begin{split}
 A_1={}&-C_3-C_4+C_8+2C_{14}+C_{15}+C_{19}+C_{21}+C_{23},\\
 A_2={}&-C_1-C_2+C_3+C_8-C_9+2C_{10}+C_{15}+C_{19}
          -C_{20}+2C_{22}+C_{24}.
\end{split}
\]
Here the $C_i$ are the global closed cusp divisors in the order
specified by their $65$ modular-unit valuations in
\path{short_cusp_divisors.m}. All cusps occurring in these
two expressions are $F$-rational. For each $i$, we have
$h^0(X,A_i)\leq1$.
\end{lemma}
\begin{proof}
Let $c_0$ be the first rational cusp, and put
$\mathcal C=\langle[C_i-(\deg C_i)c_0]:1\leq i\leq24\rangle$.
Its order is $20253807$. We use the established rank-zero
input for $J(F)$ and injectivity of torsion at the two split
odd primes $7$ and $13$.

There are four remaining normalized pairs of reductions. Their labels
at $13$ are $106,1168,1202,1520$, and the two possible source
patterns at $7$ are $3514$ and $3561$.
Both source classes $[D_7-5c_0]$ belong to the reduction of
$\mathcal C$. Since rational torsion injects at $7$, this forces
the global class $[D-5c_0]$ itself to belong to $\mathcal C$.
The classes labelled $1168$ and $1520$ at $13$ do not belong
to its reduction, so these two cases are impossible.

For labels $106$ and $1202$, the actual unmultiplied classes
at both primes agree with $[A_1-5c_0]$ and $[A_2-5c_0]$,
respectively. The normalization uses the actual cusp
permutation and the correction $5(gc_0-c_0)$.
The lattice calculations retain the full native class
coordinates; no cuspidal generation statement for $J(F)$
is assumed. Injectivity at $7$ now gives the claimed global
linear equivalences.

Finally, \cite[Table~1]{NV} gives
$\operatorname{gon}_{\mathbb F_7}X\geq6$, hence
$\operatorname{gon}_F X\geq6$. It follows that
$h^0(X,A_i)\leq1$.
\end{proof}


\begin{proposition}\label{prop:decic318-classification}
Put $F=\mathbb Q(\zeta_3)$ and $X=X_1(3,18)_F$.
The effective $F$-rational divisors of degree $5$ on $X$
supported away from the cusps form exactly two orbits under
\[
 G=\{(P,Q)\mapsto(aP+6bQ,cP+dQ):
 b,c\in\mathbb Z/3\mathbb Z,\ d\in\{1,5,7\},\ a=d^{-1}\pmod3\}.
\]
Every such divisor is a closed quintic point.
The two orbits are represented by the following curve over
$L=F(u)$ and its image under the involution $\tau$ of
Lemma~\ref{lem:decic-318-quartic-fibre}:
\[
\begin{split}
 0={}&3u^5+(1+5\zeta_3)u^4+(6\zeta_3-3)u^3
           -(6\zeta_3+9)u^2-(5\zeta_3+4)u+3,\\
 19t={}&(-25+4\zeta_3)u^4+(12-65\zeta_3)u^3
         +(20+31\zeta_3)u^2\\
       &+(69+30\zeta_3)u-33+3\zeta_3,\\
 E:\quad y^2+{}&\bigl((u+2)t+1-u\bigr)xy
                  +\bigl((u+1)t^2-ut\bigr)y=x^3.
\end{split}
\]
The marked points are $P_3=(0,0)$ and $Q_{18}=(-t,t^2)$.
Both represented elliptic curves have full torsion group
$C_3\times C_{18}$ over $L$.
\end{proposition}
\begin{proof}
We have $g(X)=10$ and $J_X(F)$ is finite by
\cite[Theorem~4.1]{DS}. The bound in \cite[Table~1]{NV} gives
$\operatorname{gon}_{\mathbb F_7}X\geq6$, hence
$\operatorname{gon}_F X\geq6$.
We use the two split primes $7$ and $13$.
At either prime the closed cusp signature is $1^{18},3^6$.
Both Frobenius elements generate
$\operatorname{Gal}(\mathbb Q(\zeta_{18})/F)$, so every closed
residue cusp lifts to the corresponding global cuspidal trace.
Entirely cuspidal reduction of a proposed divisor contradicts
torsion injectivity at the odd prime and the gonality bound.

The complete noncuspidal closed-place counts are
\[
\begin{array}{c|rrrrr}
 p&1&2&3&4&5\\ \hline
 7&0&27&144&513&3672\\
 13&0&54&846&7506&75384.
\end{array}
\]
They are checked by the Frobenius polynomial and by the explicit
closed-place enumerations. The group $G$ has order $27$ and acts
freely on geometric noncuspidal points. Indeed, its level changes
fix the nonzero point $9Q_{18}$ of order $2$, whereas a nontrivial
elliptic automorphism of order $3$ modulo sign cannot fix such a point.
The numbers of closed-place orbits in degrees two through five are
$(1,8,19,136)$ at $7$ and $(2,38,278,2792)$ at $13$.

We enumerate the possibilities for the noncuspidal support and
its cusp complement. Two quadratic supports may coincide.
After fixing one support by $G$, the numbers of cases are
\[
\begin{array}{c|rrrrrr|r}
 p&2+C_3&3+C_2&4+C_1&5&2+2+C_1&2+3&\text{total}\\ \hline
 7&1146&1368&342&136&486&144&3622\\
 13&2292&6498&5004&2792&1944&1692&20222.
\end{array}
\]
Here $C_j$ denotes any effective cuspidal divisor of degree $j$;
the enumeration includes both rational cusp multisets and cubic cusps.

Choose a rational cusp $c_0$, and let $\mathcal C\subset J_X(F)$
be generated by the $24$ global cusp-block classes
$[C_i-(\deg C_i)c_0]$.
Their reductions are identified by the valuations of the same
$65$ integral modular units at both primes.
The proved finite Jacobian presentations give
\[
\begin{split}
 J_X(\mathbb F_7)&\simeq C_3^3\times C_9^2\times C_{63}^2\times C_{189},\\
 \mathcal C&\simeq C_3^4\times C_{63}^3,
             \qquad |\mathcal C|=20253807.
\end{split}
\]
For $q=13,19$, Eichler--Shimura gives
$A_qJ_X(F)=0$, where $A_q=T_q-1-q\langle q\rangle$.
In the actual finite presentation at $7$, the common kernel $K$
of these two operators satisfies
\[
                       K/\overline{\mathcal C}\simeq C_3^3.
\]
The operators are computed on the cusp classes and three further
quadratic classes which together generate the full finite Jacobian.
Every relation is checked. In particular,
$3J_X(F)\subset\mathcal C$; no cuspidal generation assertion
for $J_X(F)$ is used.

The split-Hecke conditions leave $926$ cases at $7$.
Lemma~\ref{lem:decic-318-cubic-fibre} removes $684$ cubic-plus-cusp
cases, leaving $242$.
For $v=[D-5c_0]$, the pair
$(3\overline v_7,3\overline v_{13})$ must lie in the diagonal image
of $\mathcal C$. Exact integer lattice arithmetic applies this
condition to all $20222$ cases at $13$, after the cubic cases there
are removed by the same lemma. It leaves $78$ cases, with $53$
distinct noncuspidal supports.

For each of these supports $N$, the divisor
$T_7N-N-7\langle7\rangle N$ must be principal. The class of its cuspidal
complement is killed by the same operator, since it is a global
$F$-rational cusp divisor. The exact Riemann--Roch tests give
zero in $43$ of the $53$ cases.
The remaining ten supports give twelve possible cases at $13$.
Their paired reductions at $7$ are either a quartic in orbit
$12$ or $17$ together with a rational cusp, or one of two
quadratic-plus-cubic cases. The quartic cases are excluded by
Lemma~\ref{lem:decic-318-quartic-fibre}.
Lemma~\ref{lem:decic318-last-cuspidal-classes} then reduces the
remaining possibilities to the two global degree-five classes
$A_1,A_2$ displayed there. This last step uses unmultiplied
cuspidal membership and the actual cusp permutations.

Exact characteristic-zero Riemann--Roch gives
\[
                            h^0(X,A_1)=h^0(X,A_2)=1.
\]
The first section is checked independently by its finite and
infinite ideal pole inequalities. Its norm to the rational
function field has a quintic numerator factor. Specializing the
source equation and the section at this factor gives a linear
greatest common divisor in the second coordinate.
The resulting point is precisely the displayed curve over $L$.
The quintic defining $L/F$ is irreducible, and direct substitution
checks the source equation, the two marked orders, and their
primitive Weil pairing.
At $13$ its effective divisor reduces to the unique noncuspidal
quintic representing $A_1$.

The involution $\tau$ interchanges the two $G$-orbits.
Indeed, its cusp permutation gives
$\tau A_1\sim gA_2$ for an explicitly retained $g\in G$.
This equality is checked in the finite cuspidal group at $7$;
Manin--Drinfeld and odd-prime torsion injectivity give the global
equality. The direct marked-point transformation at $13$ gives
the same relation. Hence the unique effective divisor in the
second class is an automorphic image of $\tau D_1$.
The two orbits are distinct by their reductions at $13$.
Since $5$ is coprime to $27$, their stabilizers in $G$ are trivial.

Finally, the displayed points generate $C_3\times C_{18}$.
The absolute degree-ten defining polynomial is squarefree modulo
$7$, with factor degrees $2,2,3,3$.
At either quadratic residue prime the displayed curve has good
reduction and exactly $54$ points over $\mathbb F_{49}$.
The same holds for its $2$-isogenous curve $\tau E$.
Torsion injectivity therefore proves that both full torsion groups
have order $54$, as claimed.

The certificate retains the proved native class presentations,
the complete finite enumerations, and the programs and verification
logs for the principal-divisor identities used by the forward class maps.
All $2394$ cubic and $36$ quartic differential restrictions are
replayed independently. For the $43$ zero-dimensional Hecke tests,
a second native calculation reconstructs $292$ prime references;
independent arithmetic checks $45$ distinct marked sources and all
$360$ geometric $7$-kernels, including both marked images and every
residue norm. The unmultiplied class comparisons, global cusp
bindings, global section, and exact number-field reductions are
also checked independently.
\end{proof}

\begin{corollary}\label{cor:decic318-no-larger}
Over a number field of degree $10$, the groups $C_3\times C_{36}$
and $C_6\times C_{18}$ do not occur as torsion subgroups.
\end{corollary}
\begin{proof}
Either group contains $C_3\times C_{18}$.
The Weil pairing gives $F\subset K$, and the corresponding degree-five
trace divisor is covered by Proposition~\ref{prop:decic318-classification}.
Every represented elliptic curve has full torsion of order $54$.
\end{proof}
